\section{Introduction}

Gauss's work on the constructibility of regular polygons is one of his most famous achievements. At age 19, he proved that a regular heptadecagon (17-gon) is constructible with ruler and compass, and more generally characterized all constructible regular $n$-gons. This discovery was so important to Gauss that he requested a 17-gon be inscribed on his tombstone.

In this chapter, we cover:
\begin{itemize}
    \item Roots of unity and cyclotomic polynomials
    \item Gauss's criterion for constructible polygons
    \item The 17-gon construction using Gaussian periods
    \item General $n$-gon construction when possible
    \item Implementation of the construction
\end{itemize}

\section{Roots of Unity}

The vertices of a regular $n$-gon inscribed in the unit circle are the $n$-th roots of unity:
\[
\zeta_n^k = e^{2\pi i k / n}, \quad k = 0, 1, \ldots, n-1.
\]

\begin{definition}[Cyclotomic Polynomial]
The $n$-th cyclotomic polynomial is
\[
\Phi_n(x) = \prod_{\substack{1 \le k \le n \\ \gcd(k,n)=1}} (x - \zeta_n^k).
\]
It is the minimal polynomial of $\zeta_n$ over $\Q$, has degree $\varphi(n)$, and is irreducible.
\end{definition}

\section{Gauss's Criterion}

\begin{theorem}[Gauss-Wantzel Theorem]
A regular $n$-gon is constructible with ruler and compass if and only if
\[
n = 2^k p_1 p_2 \cdots p_m
\]
where $p_i$ are distinct Fermat primes: $F_k = 2^{2^k} + 1$.
\end{theorem}

Known Fermat primes: $3, 5, 17, 257, 65537$. The next $2^{2^5}+1 = 4294967297 = 641 \times 6700417$ is composite.

\section{The 17-gon Construction}

Gauss used Gaussian periods to express $\cos(2\pi/17)$ in nested square roots. Let $\zeta = e^{2\pi i/17}$. The periods of length 2 are:
\begin{align*}
\eta_1 &= \zeta + \zeta^9 + \zeta^{13} + \zeta^{15} + \zeta^{16} + \zeta^8 + \zeta^4 + \zeta^2 \\
\eta_2 &= \zeta^3 + \zeta^{10} + \zeta^5 + \zeta^{11} + \zeta^{14} + \zeta^7 + \zeta^{12} + \zeta^6
\end{align*}

These satisfy $\eta_1 + \eta_2 = -1$, $\eta_1 \eta_2 = -4$. Then $\eta_1 = \frac{-1 + \sqrt{17}}{2}$, $\eta_2 = \frac{-1 - \sqrt{17}}{2}$.

Continuing this process yields the nested square root expression for $\cos(2\pi/17)$.

\section{Python Implementation}

In \texttt{python/heptadecagon.py}:

\begin{lstlisting}[style=python, caption={Key functions in python/heptadecagon.py}]
cyclotomic_poly(n)              # Phi_n(x) coefficients
constructible_ngons(limit)      # List constructible n
heptadecagon_vertices()         # Coordinates of 17-gon
gaussian_periods(n, length)     # Gaussian periods
cos_2pi_over_17()              # Exact expression
construct_polygon(n)            # General construction
\end{lstlisting}

\section{Numerical Examples}

\begin{example}[17-gon Vertices]
\begin{lstlisting}[style=python]
>>> from heptadecagon import heptadecagon_vertices
>>> verts = heptadecagon_vertices()
>>> for i, (x, y) in enumerate(verts[:4]):
...     print(f"V{i}: ({x:.6f}, {y:.6f})")
V0: (1.000000, 0.000000)
V1: (0.932472, 0.361242)
V2: (0.739009, 0.673696)
V3: (0.445738, 0.895163)
\end{lstlisting}
\end{example}

\begin{example}[Constructible $n$-gons]
\begin{lstlisting}[style=python]
>>> from heptadecagon import constructible_ngons
>>> constructible_ngons(50)
[3, 4, 5, 6, 8, 10, 12, 15, 16, 17, 20, 24, 30, 32, 34, 40, 48]
\end{lstlisting}
\end{example}

\section{Exercises}

\begin{exercise}
Derive the nested square root expression for $\cos(2\pi/17)$ step by step using Gaussian periods.
\end{exercise}

\begin{exercise}
Implement the construction of the 257-gon using Gaussian periods. How many square root nestings are needed?
\end{exercise}

\begin{exercise}
Prove that $\Phi_n(x)$ is irreducible over $\Q$ using Eisenstein's criterion on $\Phi_p(x+1)$ for prime $p$.
\end{exercise}

\begin{exercise}
Show that the coordinates of the 17-gon vertices lie in the field $\Q(\sqrt{17}, \sqrt{34-2\sqrt{17}}, \ldots)$.
\end{exercise}

\section{References}

\begin{itemize}
    \item \cite{gauss-disq} -- \textit{Disquisitiones Arithmeticae}, Articles 355--366
    \item \cite{cox-galois} -- Cox, \textit{Galois Theory}, Ch. 11
    \item \cite{stein-17gon} -- Stein, \textit{The 17-gon: A Study in Constructibility}
    \item \cite{wantzel} -- Wantzel (1837), proof of necessity
\end{itemize}